% Generated draft: requires mathematical and source review.
\documentclass[11pt]{article}
\usepackage[margin=25mm]{geometry}
\usepackage{fontspec}
\setmainfont{Latin Modern Roman}
\usepackage{amsmath,amssymb,mathtools,unicode-math}
\setmathfont{Latin Modern Math}
\usepackage{enumitem,needspace,xurl,hyperref}
\setlength{\parindent}{0pt}
\setlength{\parskip}{5pt}
\setlength{\emergencystretch}{3em}
\newcommand{\N}{\mathbb{N}}
\newcommand{\Z}{\mathbb{Z}}
\newcommand{\Q}{\mathbb{Q}}
\newcommand{\R}{\mathbb{R}}
\newcommand{\C}{\mathbb{C}}
\newcommand{\F}{\mathbb{F}}
\newcommand{\A}{\mathbb{A}}
\newcommand{\PP}{\mathbb{P}}
\newcommand{\E}{\mathbb{E}}
\newcommand{\eps}{\varepsilon}
\newcommand{\dd}{\,\mathrm{d}}
\begin{document}
\section*{UnsolvedMath ID: 30005018}
\textbf{Source problem code:} catalogue entry
\subsection{Definitions and mathematical statement}
Fix an integer \(n\geq 3\). Let \(P\subset\R^2\) be a polygon with cyclically indexed vertices \(v_1,\ldots,v_n\). The source does not specify the admissible class of polygons, including convexity, simplicity, or nondegeneracy conditions. The following formulation is therefore conditional on specifying that class and the coordinate conventions.

For each admissible \(P\), let \(D_P\subseteq P\) be the domain on which its mean-value barycentric coordinates \(\beta_1,\ldots,\beta_n\) are defined, including any prescribed boundary extensions. These are the mean-value coordinate functions named in the source, normalized by
\[
\sum_{i=1}^n\beta_i(x)=1,\qquad
x=\sum_{i=1}^n\beta_i(x)v_i\quad(x\in D_P).
\]
These identities alone do not define mean-value coordinates uniquely; their defining convention and boundary values must also be supplied. Define
\[
\Phi_P:D_P\longrightarrow\mathbb P^{n-1}(\R),\qquad
x\longmapsto[\beta_1(x):\cdots:\beta_n(x)].
\]
Normalization ensures that these projective coordinates do not vanish simultaneously. For a specified choice \(K\in\{\R,\C\}\), regard this image as a subset of \(\mathbb P^{n-1}(K)\), and define its homogeneous vanishing ideal and projective Zariski closure by
\[
I_{P,K}=\bigoplus_{d\geq0}
\left\{F\in K[X_1,\ldots,X_n]_d:
F(\beta_1(x),\ldots,\beta_n(x))=0
\text{ for every }x\in D_P\right\},
\qquad
Y_{P,K}=\operatorname{Proj}\bigl(K[X_1,\ldots,X_n]/I_{P,K}\bigr).
\]
The first problem is to determine homogeneous equations defining \(Y_{P,K}\), or, in the stronger ideal-theoretic interpretation of the source's request for algebraic relations, to determine generators of \(I_{P,K}\), for arbitrary admissible \(n\)-gons \(P\). The source states that mean-value coordinates are algebraic functions despite descriptions using expressions such as \(\tan(\theta/2)\). It does not specify their algebraic branches, the ground field for the closure, or whether equations are requested set-theoretically or ideal-theoretically.

For the second problem, let \(w_1,\ldots,w_n\) denote the Wachspress barycentric coordinates of an admissible polygon \(P\), on their domain of definition, and let
\[
S_{P,K}=\overline{\left\{[w_1(x):\cdots:w_n(x)]:x\text{ lies in that domain}\right\}}^{\,\mathrm{Zar}}
\subseteq\mathbb P^{n-1}_K
\]
be its Wachspress surface, with the reduced scheme structure for the starting surface. The source does not supply the coordinate formula or the precise admissibility assumptions for this construction. It reports, attributing the result to Irving and Schenck, that Wachspress surfaces of \(n\)-gons have a common Hilbert polynomial; denote that reported polynomial by \(H_n\) on the intended class of polygons. Its formula and the hypotheses of the attributed result are not supplied.

A scheme-theoretic flat limit can be expressed as a special fiber \(\mathcal S_0\) of a closed subscheme
\[
\mathcal S\subseteq\mathbb P^{n-1}_A
\]
flat over a discrete valuation \(K\)-algebra \(A\), whose generic fiber belongs to the chosen family of Wachspress surfaces. Here \(\mathcal S_0=\mathcal S\times_A A/\mathfrak m_A\), and \(\mathfrak m_A\) is the maximal ideal of \(A\). With Hilbert polynomial fixed to \(H_n\), the second problem is to describe the flat limits arising as polygons degenerate and to construct a natural, meaningful compactification of the corresponding parameter or moduli space, explaining its boundary objects. The source does not specify the allowed degenerations, the equivalence relation on polygons or surfaces, or a precise moduli functor. Consequently it does not determine a unique compactification problem. Limiting objects must be treated as schemes, since the source explicitly asks for flat limits.

The opening question about higher dimensions does not specify a higher-dimensional analogue, its hypotheses, or its intended target. The two explicit problems above concern polygons in the plane.
\subsection{Short English statement}
For each admissible planar \(n\)-gon, find the homogeneous algebraic relations among its mean-value barycentric coordinates. Also describe the scheme-theoretic flat limits of Wachspress surfaces as polygons degenerate, and formulate and construct a meaningful compactified parameter or moduli space. Polygon hypotheses, coordinate conventions, ground field, allowable degenerations, and moduli equivalences remain unspecified in the supplied source.
\subsection{Sources}
\begin{itemize}
\item Supplied local UnsolvedMath registry record, permanent ID \(30005018\), title \emph{Mean-Value Coordinate Equations and Wachspress-Surface Moduli}; Oberwolfach Reports Open Problems collection; exact source problem code \texttt{catalogue entry}. The registry supplies the status \texttt{partially\_solved} and statement status \texttt{corrected\_verified}. Its statement-verification field reports comparison with the official report; this draft does not independently verify that comparison.
\item Supplied report citation: \emph{Mini-Workshop: Interpolation, Approximation, and Algebra (2023)}, Oberwolfach Reports, DOI 10.4171/owr/2022/7, pp.\ 383--428. \url{https://doi.org/10.4171/owr/2022/7}. The record gives proposed year 2022; the citation's 2023 date and the DOI's 2022 component are retained without reconciliation. The background describes Oberwolfach Reports as documenting workshops and their open-problem discussions.
\item Attribution in the supplied original statement: both questions are attributed to Sottile; the mean-value question is described as inspired by Schafer's talk and by Corey Irving's Ph.D. work on analogous questions for Wachspress coordinates. The common-Hilbert-polynomial assertion is attributed to Irving and Schenck. The report also says Irving and Sottile studied examples of limits ``a decade ago,'' without a conclusive outcome; this is historical source wording, not a current literature finding.
\item Supplied literature source: Kohn--Ranestad, projective geometry of Wachspress coordinates. \url{https://arxiv.org/abs/1904.02123}.
\item Supplied literature source: Kohn--Ranestad, Wachspress-coordinate map. \url{https://doi.org/10.1007/s10208-019-09441-z}.
\item Literature assessment checked at \texttt{2026-08-22}: supplied dated curation, not independently verified. It labels the problem partially solved and states that projective and adjoint theory describes Wachspress maps and varieties in broad generality and provides substantial three-dimensional results, while equations and a complete description in arbitrary higher dimensions remain open. The background identifies this as dated literature triage derived from the supplied Oberwolfach curation and requests independent verification of time-sensitive status and cited claims. Its implications for the two precise planar questions are not established by the supplied record.
\end{itemize}
\end{document}
